\scanpage{45}
\mspage{41}
% archive sheet number: 44

Supposons maintenant qu'on ait
\[
X\ll Y\qquad (X,Y\in\mathcal A)
\]
on a alors une application croissante canonique
\[
\unsure{\varphi|\phi}:\widetilde X\longrightarrow\widetilde Y
\]
telle que, pour \deleted{tout} $X'\in\widetilde X$, on ait
\[
X'\unsure{\overset{\circ}{\ll}|\overset{\circ}{\Subset}}\deleted{\varphi(X')\overset{\mathrm{def}}{=}}Y'\overset{\mathrm{def}}{=}\varphi(X').
\]
En termes de \unsure{déployements}, on peut interpréter $\varphi$ par une application croissante $\unsure{\psi|\tau}$
\[
\psi:\mathrm{D\acute{e}p}(X)\longrightarrow\mathrm{D\acute{e}p}(Y)
\]
satisfaisant
\[
\forall S\in\mathrm{D\acute{e}p}\,X,\qquad S\leq\psi(S)
\]
(ici on a noté $\leq$ la relation $\subset$ dans $\Sigma$).

Cette application $\psi$ est uniquement déterminée par cette propriété, car \deleted{si} \deleted{\unsure{peut avoir|pour avoir}}
\[
S\leq T,T',\qquad\text{avec }T,T'\in\mathrm{Dep}(Y)
\]
on a $T=T'$ (et ceci, pour peu seulement que $S$ soit un élément de $\Sigma$, $\ne\varnothing_\Sigma$). En effet, si $T\ne T'$, on a $T\mathrel{|\circ|}T'$,
